\section{May 26}

\subsection{Motivic cohomology}

  \paragraph{Idea:} Let $X$ be a topological space, $C_\bullet(X)$ be the singular chain complex of $X$. 
  Then $H^n(X,\bbZ) := H^n(\Hom(C_\bullet(X),\bbZ))$.
  Let $\sfD(\bbZ)$ be the derived category of $\bbZ$, i.e., the category of chain complexes of abelian groups with quasi-isomorphisms inverted. 
  Then
  \[ H^n(X,\bbZ) = \Hom_{\sfD(\bbZ)}(C_\bullet(X),\bbZ[n]) \]
  and 
  \[ \bbZ[n] = C_\bullet(S^n)/C_\bullet(\mathrm{pt}) = C_{\bullet}(\bigwedge^n S^1)/C_\bullet(\mathrm{pt}) = (C_\bullet(S^1)/C_\bullet(\mathrm{pt}))^{\otimes n}. \]
  In summary, we have
  \[ H^n(X,\bbZ) = \Hom_{\sfD(\bbZ)}(C_\bullet(X), (C_\bullet(S^1)/C_\bullet(\mathrm{pt}))^{\otimes n}). \]
  Voevodsky's idea is to construct in AG analogues of $C_\bullet(X), \sfD(\bbZ), H^n(X,\bbZ)$, 
  which will be motives $\sfM(X)$, the triangulated category $DM_{Nis}^{eff,-}(k,\bbZ)$ of motives, motivic cohomology $H^{p,q}(X,\bbZ)$ so that above formula still holds.


\subsection{Finite correspondences and presheaves with transfers}

  Let $\kk$ be a field, $\Sm_k$ be the category of smooth schemes of finite type over $k$.

  \begin{definition}
    Let $X,Y\in \Sm_k$. 
    (When $X$ is connected, )
    An elementary correspondence from $X$ to $Y$ is an integral closed subvariety $Z\subset X\times Y$ such that the projection $Z\to X$ is finite and surjective.
    (When $X$ is not connected, an elementary correspondence from $X$ to $Y$ is an elementary correspondence from a connected component of $X$ to $Y$.)

    Define $\Cone(X,Y)$ to be the free abelian group generated by elementary correspondences from $X$ to $Y$. 
    The elements of $\Cone(X,Y)$ are called finite correspondences from $X$ to $Y$.
  \end{definition}

  \begin{example}
    Let $f: X\to Y$ be a morphism in $\Sm_k$. 
    Then the graph $\Gamma_f\subset X\times Y$ of $f$ is an elementary correspondence from $X$ to $Y$.
    If $f$ is finite and surjective, then $\Gamma_f^\mathrm{t}$ is also an elementary correspondence from $Y$ to $X$, where $\Gamma_f^\mathrm{t}$ is the transpose of $\Gamma_f$.
  \end{example}

  \begin{example}
    Let $X,Y \in \Sm_k$ and $Z\subset X\times Y$ be a closed subscheme that is finite and surjective over $X$. 
    Then $[Z] \in \Cone(X,Y)$ is a finite correspondence from $X$ to $Y$.
  \end{example}

  \begin{lemma}
    Let $X,Y,Z\in \Sm_k$, $\alpha\in \Cone(X,Y)$, $\beta\in \Cone(Y,Z)$. 
    Then 
    \[ \beta\circ \alpha := p_{XZ*}(p_{XY}^*\alpha \cdot p_{YZ}^*\beta) \in \Cone(X,Z), \]
    where $p_{XY}, p_{YZ}, p_{XZ}$ are the projections from $X\times Y\times Z$ to $X\times Y, Y\times Z, X\times Z$ respectively.
  \end{lemma}
  \begin{proof}
    We may assume $\alpha = [W]$ and $\beta = [V]$ for some elementary correspondences $W\subset X\times Y$ and $V\subset Y\times Z$. 
    As for $V \times_Y W \subset X\times Y\times Z$, we have the following commutative diagram:
    \[ \begin{tikzcd}
      V \times_Y W \ar[r] \ar[d, "\tilde{g}"'] & W \ar[d, "g"] \ar[r] & Z \\
      V \ar[r] \ar[d, "f"'] & Y &  \\
      X & &
    \end{tikzcd} \]
    Since $f,g$ are finite and surjective, so is $f\circ \tilde{g}$.
    Then $p_{XZ*}(V \times_Y W)$ is a closed subscheme of $X\times Z$ that is finite and surjective over $X$, which means $[p_{XZ*}(V \times_Y W)] \in \Cone(X,Z)$.
  \end{proof}

  We will denote $\Cor_k$ the category whose objects are those of $\Sm_k$ and morphisms from $X$ to $Y$ are $\Cone(X,Y)$, with composition defined as above.

  \begin{proposition}
    We have the following properties of $\Cor_k$:
    \begin{enumerate}
      \item $\forall X\in \Sm_k$, the identity morphism of $X$ in $\Cor_k$ is $[\Delta_X]$, where $\Delta_X\subset X\times X$ is the diagonal.
      \item $\Sm_k \to \Cor_k$ sending $X$ to $X$ and $f: X\to Y$ to $[\Gamma_f]$ is a faithful functor.
      \item $\Cor_k$ is a symmetric monoidal category with monoidal structure given by the product of schemes, i.e., $X\otimes Y := X\times Y$ for $X,Y\in \Cor_k$.
    \end{enumerate}
  \end{proposition}

  \begin{example}
    Let $X \in \Sm_k$.
    Then 
    \begin{itemize}
      \item $\Cor(X,\Spec \kk) = \bigoplus_{i} \bbZ [X_i]$, where $X_i$ are the connected components of $X$.
      Then structure morphism $X\to \Spec \kk$ corresponds to $\sum_i [X_i] \in \Cor(X,\Spec \kk)$.

      \item $\Cor(\Spec \kk,X) = \bigoplus_{x\in X} \bbZ [x]$, where $x$ runs through the closed points of $X$.
    \end{itemize}
    Let $x \in X$ be a closed point.
    \begin{itemize}
      \item $\Spec \kk \xrightarrow{x} X \to \Spec \kk$ corresponds to $[\kappa(x):\kk] \in \Cor(\Spec \kk,\Spec \kk)$, where $\kappa(x)$ is the residue field of $x$.
      \item $X \to \Spec \kk \xrightarrow{x} X$ corresponds to $[X \times_{\Spec \kk} \{x\}] \in \Cor(X,X)$.
    \end{itemize}
  \end{example}

  \begin{definition}
    A presheaf with transfers is an additive contravariant functor $F: \Cor_k \to \Ab$.
    Define $\PST(\kk)$ to be the category of presheaves with transfers.
  \end{definition}

  By additivity,  we get: $\forall X,Y\in \Sm_k, \Cone(X,Y) \to \Hom_{\Ab}(F(Y),F(X))$ sending $\alpha$ to $F(\alpha)$ is a group homomorphism.
  Hence there is a pairing
  \[ \Cone(X,Y) \times F(Y) \to F(X),\quad (\alpha,a) \mapsto F(\alpha)(a). \]

  \begin{lemma}
    $\PST(\kk)$ is an abelian category with enough projectives and injectives.
    In particular, (co)homological groups of chain complexes of presheaves with transfers are well-defined.
  \end{lemma}
  \begin{proof}
    See \cite[Section 2.1]{MVW}.
  \end{proof}

  \begin{example}
    Let $A$ be a constant presheaf on $\Sm_k$.
    Then $A$ is a presheaf with transfers by defining 
    \[ A(X) := A, \quad F([V]) := \deg(V/X) \cdot \mathrm{id}_A, \]
    where $V\subset X\times Y$ is an elementary correspondence from $X$ to $Y$ and $\deg(V/X)$ is the degree of the finite morphism $V\to X$.
  \end{example}

  \begin{example}
    The sheaves $\calO^*, \calO \in \PST(\kk)$ by defining $\calO^*(X) := \Gamma(X,\calO_X)^*$, $\calO(X) := \Gamma(X,\calO_X)$ and 
    \begin{align*}
      \calO^*([V]) & := \bigl( \calO^*(Y) \to \calO^*(W) \xrightarrow{\Nm_{V/X}} \calO^*(X) \bigr), \\
      \calO([V]) & := \bigl( \calO(Y) \to \calO(W) \xrightarrow{\Tr_{V/X}} \calO(X) \bigr).
    \end{align*}
    (Here we use the fact that $X$ is smooth and hence normal.)
  \end{example}

  \begin{example}
    For any $i$, $\CH^i(-) \in \PST(\kk)$ by defining $\CH^i(X) := \CH^i(X)$ and 
    \[ \CH^i(\alpha) := \bigl( \CH^i(Y) \xrightarrow{p_Y^*} \CH^i(W) \xrightarrow{p_{X*}(\alpha \cdot -)} \CH^i(X) \bigr). \]
    (Note that here $X/\kk$ does not have to be proper since the pushforward $p_{X*}$ is finite on $\alpha$'s support and hence well-defined.)
  \end{example}

  More examples are provided by representable functors.

  \begin{definition}
    Let $X\in \Sm_k$.
    Define $\bbZ_{tr}(X) \in \PST(\kk)$ by setting $\bbZ_{tr}(X)(Y) := \Cor(Y,X)$ for $Y \in \Sm_k$ and $\bbZ_{tr}(X)(\alpha) := \alpha \circ -$ for $\alpha \in \Cor(Y,Z)$.
  \end{definition}

  \begin{proposition}
    We have the following properties of $\bbZ_{tr}(X)$:
    \begin{enumerate}
      \item $\bbZ_{tr}(X)$ is a projective object in $\PST(\kk)$ (by Yoneda's lemma: $\Hom_{\PST(\kk)}(\bbZ_{tr}(X),F) \cong F(X)$ for $F\in \PST(\kk)$).
      \item $\forall X,Y\in \Sm_k, \forall f \in \Cone(X,Y)$, we have $f_* : \bbZ_{tr}(X)(U) \to \bbZ_{tr}(Y)(U)$ defined by $f_*(\alpha) := f \circ \alpha$.
    \end{enumerate}
  \end{proposition}

  \begin{example}
    $\bbZ_{tr}(\Spec \kk) = \bbZ$ is the constant presheaf with transfers.
  \end{example}

  \begin{definition}
    Let $(X,x)$ be a pointed smooth scheme, i.e., $X\in \Sm_k$ and $x\in X(\kk)$ is a $\kk$-point.
    Define $\bbZ_{tr}(X,x) := \coker(\bbZ_{tr}(\Spec \kk) \xrightarrow{x_*} \bbZ_{tr}(X))$.
  \end{definition}

  The push-forward along the structure morphism $X\to \Spec \kk$ induces a retraction of $x_*$, which means $\bbZ_{tr}(X) \cong \bbZ_{tr}(X,x) \oplus \bbZ$.

  Let $(X_1,x_1), \cdots, (X_n,x_n)$ be pointed smooth schemes.
  Define 
  \[ \bbZ_{tr}((X_1,x_1) \wedge \cdots \wedge (X_n,x_n)) := \coker\left( \bigoplus_{i=1}^n \bbZ_{tr}(X_1\times \cdots \widehat{X_i} \times \cdots \times X_n) \to \bbZ_{tr}(X_1\times \cdots \times X_n) \right). \]
  Then $\bbZ_{tr}((X_1,x_1) \wedge \cdots \wedge (X_n,x_n))$ is a direct summand of $\bbZ_{tr}(X_1\times \cdots \times X_n)$ and hence projective in $\PST(\kk)$.

  By convention, we will write $\bbZ_{tr}((X,x)^{\wedge n}) = 0$ if $n < 0$.
  Set $\bbG_m := (\bbA^1\setminus \{0\}, 1)$.
  $\bbZ_{tr}(\bbG_m^{\wedge n})$ will play an important role in the construction of motives.

\subsection{Simplicial structures}

  Let $\Delta$ be the simplicial category, i.e., the category whose objects are $[n] := \{0,1,\cdots,n\}$ for $n\geq 0$ and morphisms from $[m]$ to $[n]$ are order-preserving maps from $\{0,1,\cdots,m\}$ to $\{0,1,\cdots,n\}$.

  \begin{definition}
    A \emph{simplicial object} (resp. \emph{cosimplicial object}) in a category $\bfC$ is a contravariant (resp. covariant) functor $\Delta \to \bfC$.
  \end{definition}

  A simplicial object $F: \Delta \to \bfC$ has the following data:
  \begin{itemize}
    \item face maps $\partial_i: F([n]) \to F([n-1])$ for $0\leq i \leq n$ corresponding to the morphism $[n-1] \to [n]$ that misses $i$;
    \item degeneracy maps $\sigma_i: F([n]) \to F([n+1])$ for $0\leq i \leq n$ corresponding to the morphism $[n+1] \to [n]$ that hits $i$ twice.
  \end{itemize}

  Recall $\Delta^n := \Spec \kk[t_0,\cdots,t_n]/(t_0+\cdots+t_n-1)$ is the algebraic $n$-simplex.
  Then $\Delta^\bullet: \Delta \to \Sm_k$ sending $[n]$ to $\Delta^n$ is a cosimplicial object in $\Sm_k$.
  Let $F \in \PST(\kk)$.
  Then $\forall U \in \Sm_k$, $F(U\times \Delta^\bullet)$ is a simplicial abelian group, $C_\bullet F: U \mapsto F(U\times \Delta^\bullet)$ is a simplicial presheaf with transfers.
  We get $C_*F: \Cor^{\op}_k \to \sfC(\Ab)$ sending $U$ to the Moore complex of $F(U\times \Delta^\bullet)$:
  \[ \cdots \to F(U\times \Delta^2) \xrightarrow{d_2} F(U\times \Delta^1) \xrightarrow{d_1} F(U\times \Delta^0), \]
  where $d_n := \sum_{i=0}^n (-1)^i F(\partial_i)$.

  In this case, $\forall n$, $H_n(C_*F)$ is a presheaf with transfers.

  The motive $\sfM(X)$ of $X\in \Sm_k$ will be defined as the class of $C_*\bbZ_{tr}(X)$ in a certain triangulated category $DM_{Nis}^{eff,-}(k,\bbZ)$ constructed from the derived category of $\PST(\kk)$.
  Moreover, define $\bbZ(n) := C_*\bbZ_{tr}(\bbG_m^{\wedge n})[-n]$.
